\section{Conclusion}
\label{sec:Con}
\par Results found a moderate positive correlation of 0.5286 which is evidenced that AI can be used in a personality assesment. However, ZSC is shown that is suitably used as a support tool rather than a complete human recruiter replacement. Reduction on time when scoring open ended responses based on a specific Big-Five trait is massive. AI scoring was executed under a second when evaluating user's open ended response, while a human recruiter took 1-2 minutes to evaluate and score. Despite the Pearson Correlation score, it was also found AI's inability to identify linguistic and industry-speicifc terms which influenced AI's scoring. For instance, AI failed to recognise terms such as "SLA" and "CRM" in one of the user's repsonse. 

\par A primary limitation identified was small sample size which consisted of 4 users. Due to the sample size, this might have changed the validation of the prototype's use in the persoanlity assesment in the context of job recruitment. As previously stated, ZSC lacks understading in industrial specifcs due to model's general training sources. 1 human recruiter resulted in possible human bias which impacted human recruiter's scoring. In addition. BART-MNLI was only comapred in context of AI scoring which negatively impacts validation of AI's role in personality assement. 

\par Models should be fine tuned in order to improve identifying speicifc industrial terms such as HR-specific datasets. Alternatively, industry-specifc models should be used for personality assesment to combat techincal jargon. Finally, sample size should be increased by inviting more participants or make use of datasets that contains user repsonses used for personality assesment.